\documentclass[12pt,a4paper]{report}

\usepackage[utf8]{inputenc}
\usepackage{geometry}
\geometry{top=1in, bottom=1in, left=1in, right=1in}
\usepackage{float}
\usepackage{array}
\usepackage{graphicx}
\usepackage{svg}
\usepackage{booktabs} % إضافة مكتبة الجداول الاحترافية

\begin{document}
\tableofcontents
\setcounter{chapter}{5}
\chapter{Mobile Application}

\section{App Plan}

The development of the LUMO mobile application commenced with a comprehensive analysis of the primary stakeholders: parents and healthcare professionals. By identifying the specific operational, emotional, and logistical challenges faced by these users, the development team was able to design targeted, technology-driven solutions. The resulting application serves as a central hub, addressing these core problems through advanced software features.

\subsection{Parent Perspective}

Parents of autistic children require continuous support, clear insights into their child's progress, and reliable tools to manage therapy. Table \ref{tab:parent_perspective} outlines the identified problems and the corresponding solutions implemented in the LUMO application to assist parents.

\begin{table}[H]
    \centering
    \renewcommand{\arraystretch}{1.5}
    \caption{Problem vs. Solution Matrix: Parent Perspective}
    \label{tab:parent_perspective}
    \begin{tabular}{|p{0.35\textwidth}|p{0.55\textwidth}|}
        \hline
        \textbf{Identified Problem} & \textbf{LUMO Solution} \\
        \hline
        Difficulty in communicating with the doctor and other parents. & A Global Community feed (Posts, Likes, Comments) and a Real-time Chat system. \\
        \hline
        Struggle to understand the child's progress over time. & An AI-Powered Dashboard that visualizes Focus Scores, Eye Contact (Gaze tracking), and Emotions from completed sessions. \\
        \hline
        Not knowing how to deal with sudden behavioral changes. & LUMO's intelligent RAG-based Chatbot with memory, providing instant medical tips and behavioral guidance. \\
        \hline
        Lack of visibility into how therapy sessions are structured and scheduled. & Direct access to the session schedule configured by the linked doctor, along with clear post-session summaries. \\
        \hline
    \end{tabular}
\end{table}

\newpage
\subsection{Doctor Perspective}

Healthcare professionals face the challenge of managing multiple patients while requiring precise data to tailor therapies. Table \ref{tab:doctor_perspective} details the functional limitations faced by doctors and the analytical solutions provided by the LUMO platform.

\begin{table}[H]
    \centering
    \renewcommand{\arraystretch}{1.5}
    \caption{Problem vs. Solution Matrix: Doctor Perspective}
    \label{tab:doctor_perspective}
    \begin{tabular}{|p{0.35\textwidth}|p{0.55\textwidth}|}
        \hline
        \textbf{Identified Problem} & \textbf{LUMO Solution} \\
        \hline
        Inefficiency in following up with multiple patients efficiently. & A centralized Patient Directory with a Request System (Accept/Reject links) to securely link with parents. \\
        \hline
        Difficulty in analyzing the child's response to therapy. & Access to detailed Patient Insights, showing deep analytics for each session segment (Emotions, Focus, Gaze) once processed by the AI Cloud Servers. \\
        \hline
        Inability to remotely configure and schedule therapy sessions for the robot. & A Session Builder that allows doctors to construct a ``Session Sequence'' (e.g., Study, Games, Stories, Drawing) with specific durations for each segment, which the robot later fetches and executes autonomously. \\
        \hline
    \end{tabular}
\end{table}

Following the identification of problems and the conceptualization of solutions, the second phase of planning involved the development of comprehensive software diagrams and screen maps. These visual representations were crucial in outlining the application's logical flow and structural architecture, ensuring a seamless navigation experience between the distinct user roles.

Subsequently, the UI and UX design phase was initiated. The interface design was executed using Figma, a collaborative interface design tool, to establish a modern, accessible, and medical-grade user experience. The design system heavily incorporated transparent glassmorphic elements, providing a clean and calming aesthetic suitable for the application's sensitive medical context.

\begin{figure}[H]
    \centering
    \includesvg[width=0.15\textwidth]{placeholder_figma_logo.svg}
    \caption{Figma collaborative interface design tool}
    \label{fig:figma_logo}
\end{figure}

Following the successful completion of the planning and UI/UX design phases, the next crucial step involved selecting the appropriate technological stack and defining the software architecture to bring these designs to life. The detailed technical implementation and the frameworks utilized are discussed in the subsequent sections.


\section{System Architecture}

The LUMO mobile application employs a Layered Clean Architecture combined with the Model-View-ViewModel (MVVM) design pattern. This architectural choice ensures a robust separation of concerns, scalability, and enhanced maintainability. The system is structurally divided into four primary layers, interconnected sequentially to process data and user interactions.

\begin{figure}[H]
    \centering
    \includegraphics[width=0.95\textwidth]{Chapter_6/lumo-arch.png}
    \caption{LUMO Mobile Application Architecture Diagram}
    \label{fig:app_architecture}
\end{figure}

\subsection{Layer 1: UI Layer (View)}
The presentation layer is fully decoupled from the underlying business logic. It consists of declarative Flutter widgets that observe state changes and dispatch user intents. To accommodate the specific needs of autism therapy, this layer heavily integrates customized visual components, including \texttt{fl\_chart} for rendering complex analytical visualizations (e.g., focus scores and emotion distribution) and a unified glassmorphic design system to maintain a calming, distraction-free aesthetic.

\subsection{Layer 2: State Management Layer (ViewModel)}
State orchestration is handled utilizing the \texttt{Provider} package in conjunction with \texttt{ChangeNotifier}. This layer bridges the UI and the data layers. Feature-specific ViewModels (e.g., \texttt{AnalysisViewModel}, \texttt{SessionViewModel}, \texttt{ChatViewModel}) are instantiated as factories to isolate screen-specific logic, while shared providers (such as \texttt{AuthProvider} and \texttt{PatientProvider}) are maintained as singletons to manage cross-cutting application states and ensure continuous data synchronization across the platform.

\subsection{Layer 3: Dependency Injection}
To achieve complete decoupling between components, the application utilizes the \texttt{get\_it} service locator as its centralized composition root. No class instantiates its own dependencies; instead, all services, repositories, and ViewModels are injected dynamically at runtime. This strictly enforces the dependency inversion principle, ensuring that higher-level modules orchestrate logic without being tightly bound to lower-level data implementations.

\subsection{Layer 4: Data Layer (Model and Repositories)}
The lowest layer manages all data retrieval, persistence, and serialization. It utilizes the Repository Pattern to act as the single source of truth for the ViewModels. The repositories intelligently route data requests between multiple sources:
\begin{itemize}
    \item \textbf{Remote REST APIs:} Handled via the \texttt{Dio} HTTP client to communicate with the primary Laravel backend, which manages structured relational data such as authentication, profiles, the community feed, and session configuration.
    \item \textbf{Remote Real-time Streams:} Managed via Firebase to support instantaneous chat messaging, typing indicators, and push notifications.
    \item \textbf{Local Storage:} Facilitated by \texttt{Hive} and \texttt{SharedPreferences} to cache critical user data, ensuring application resilience and fallback capabilities during network interruptions.
\end{itemize}

\subsubsection*{Backend Architecture Choice: Laravel vs. Firebase}
In the LUMO mobile application, a hybrid backend architecture was adopted to maximize performance and efficiency. Laravel (REST APIs) handles structured relational data such as user authentication, profiles, doctor-patient linking, the community feed, and session configuration. Firebase, on the other hand, is integrated specifically for real-time features, namely the Chat system and Push Notifications (FCM). This separation of concerns ensures that the relational database is not overloaded with continuous real-time listeners, while still providing a seamless and highly responsive user experience for time-sensitive interactions.

% ===================================
% NEW SECTION: APP INTEGRATION
% ===================================
\newpage
\section{App Integration}

The integration phase bridges the gap between the mobile interface and the backend services. LUMO utilizes modern tools to handle high-frequency data and real-time updates.

\subsection{Framework: Flutter \& Dart}
The application is built using the Flutter 3.x framework, leveraging Dart as the primary language. This enables a unified codebase for both Android and iOS while maintaining native performance levels required for a smooth, data-rich user experience.

\begin{figure}[H]
    \centering
    \includegraphics[width=0.1\textwidth]{placeholder_flutter_dart_logos.png}
    \caption{Flutter Framework}
    \label{fig:flutter_stack}
\end{figure}

\subsection{State Management \& Local Caching}
To ensure smooth performance during heavy analytical tasks, the application integrates \texttt{Provider} for reactive UI updates and \texttt{Hive} for high-performance NoSQL local storage. This combination allows the app to function efficiently even with intermittent network connectivity.

\begin{figure}[H]
    \centering
    \includegraphics[width=0.4\textwidth]{placeholder_state_storage_logos.png}
    \caption{Integration of Provider and Hive storage}
    \label{fig:state_management}
\end{figure}

\subsection{API Communication \& Testing}
Networking is handled via the \texttt{Dio} HTTP client, implementing interceptors for secure authentication. API endpoints were rigorously tested to ensure data integrity between the mobile client and the Laravel backend.

\begin{table}[H]
    \centering
    \caption{Key API Integration Endpoints}
    \label{tab:api_endpoints}
    \begin{tabular}{ll}
        \toprule
        \textbf{Endpoint Category} & \textbf{Functionality} \\
        \midrule
        \texttt{/auth} & Secure Login, Registration, and OTP Verification \\
        \texttt{/sessions} & Creation and retrieval of doctor-configured session schedules, and the resulting analysis reports \\
        \texttt{/posts} & Post management and user interactions \\
        \texttt{/doctor/patient-request} & Doctor-Patient linking requests \\
        \bottomrule
    \end{tabular}
\end{table}

\subsection{Real-time Communication (Firebase)}
Unlike traditional polling methods, LUMO utilizes Firebase for the chat and notification features. This provides a serverless real-time experience with low latency, essential to keep parents and doctors promptly informed.

\begin{figure}[H]
    \centering
    \includegraphics[width=0.2\textwidth]{placeholder_firebase_logo.png}
    \caption{Real-time data synchronization via Firebase}
    \label{fig:firebase_sync}
\end{figure}
\subsection{Mobile App Structure}

\noindent
\begin{minipage}[t]{0.50\textwidth}
    \vspace{0pt} % Ensures alignment at the very top
    \includegraphics[width=\linewidth]{Chapter_6/mopile-app-tree.png}
\end{minipage}%
\hfill
\begin{minipage}[t]{0.55\textwidth}
    \begin{itemize}
    \item \textbf{Core Folder:} Contains the cross-cutting infrastructure of the application, including dependency injection setup (GetIt), routing configurations, network clients (Dio), application services, and the global theme definitions.
    \item \textbf{Data Folder:} Contains the data access layer, which is divided into data models, data sources (for handling Firebase, REST APIs, and local storage), and repositories that aggregate data to act as the single source of truth.
    \item \textbf{Features Folder:} Contains all the functional modules of the app (e.g., Auth, Chat, Analysis, Session). Each feature strictly follows the MVVM pattern, separating the user interface (\texttt{view}) from the business logic and state management (\texttt{view\_model}).
    \item \textbf{Shared Folder:} Contains reusable components utilized across multiple features, including common UI widgets, global state providers (such as AuthProvider and PatientProvider), and shared mixins.
\end{itemize}
\end{minipage}

% ================
\newpage
\subsection{Dealing with APIs and Sockets}

An Application Programming Interface (API) serves as a vital communication bridge, enabling distinct software applications to interact and exchange data securely without needing to understand their internal implementations. Through APIs, the application executes CRUD (Create, Read, Update, Delete) operations utilizing standard RESTful methods.

\begin{table}[H]
    \centering
    \renewcommand{\arraystretch}{1.5}
    \caption{Standard REST API Methods utilized in LUMO}
    \label{tab:rest_methods}
    \begin{tabular}{|l|p{0.75\textwidth}|}
        \hline
        \textbf{Method} & \textbf{Description} \\
        \hline
        \textbf{GET} & Retrieve information about the REST API resource. \\
        \hline
        \textbf{POST} & Create a new REST API resource. \\
        \hline
        \textbf{PUT} & Update an existing REST API resource. \\
        \hline
        \textbf{DELETE} & Delete a REST API resource or related component. \\
        \hline
    \end{tabular}
\end{table}

To effectively manage these network requests within the Flutter ecosystem, the LUMO application implements \texttt{Dio 5.4.1} as its primary HTTP client. This library provides a robust foundation for secure authentication management and structured error handling. The mobile application communicates with the Laravel backend via the production Base URL: \texttt{https://app2.clickexpress.delivery/api/}.

\vspace{10pt}
\noindent\textbf{Examples of API Methods used in the LUMO Application:}
\vspace{5pt}

\begin{table}[H]
    \centering
    \renewcommand{\arraystretch}{1.5}
    \begin{tabular}{|p{0.46\textwidth}|p{0.46\textwidth}|}
        \hline
        \multicolumn{2}{|c|}{\texttt{/api/auth/register}} \\
        \hline
        \multicolumn{2}{|c|}{\textbf{API used in Role-Based Registration for both parent and doctor}} \\
        \hline
        \multicolumn{2}{|c|}{\textbf{Method used:} POST} \\
        \hline
        \textbf{Request:} Parents and doctors send \texttt{[role, name, email, phone]}. Parents additionally send \texttt{[child\_name]}, while doctors send \texttt{[clinic\_location, doctor\_number]}. & 
        \textbf{Response:} If status code 201, the user is registered successfully. If status code 400, there is a validation error in the submitted data. \\
        \hline
    \end{tabular}
\end{table}

\begin{table}[H]
    \centering
    \renewcommand{\arraystretch}{1.5}
    \begin{tabular}{|p{0.46\textwidth}|p{0.46\textwidth}|}
        \hline
        \multicolumn{2}{|c|}{\texttt{/api/doctor/patient-request}} \\
        \hline
        \multicolumn{2}{|c|}{\textbf{API used by doctors to securely send a link request to a patient's profile}} \\
        \hline
        \multicolumn{2}{|c|}{\textbf{Method used:} POST} \\
        \hline
        \textbf{Request:} Doctor will send the unique \texttt{patient\_id} to initiate the linkage. & 
        \textbf{Response:} If status code 200, the link request is dispatched successfully to the parent for approval. \\
        \hline
    \end{tabular}
\end{table}

\begin{table}[H]
    \centering
    \renewcommand{\arraystretch}{1.5}
    \begin{tabular}{|p{0.46\textwidth}|p{0.46\textwidth}|}
        \hline
        \multicolumn{2}{|c|}{\texttt{/api/sessions}} \\
        \hline
        \multicolumn{2}{|c|}{\textbf{API used by doctors to build and schedule a complete session for the robot}} \\
        \hline
        \multicolumn{2}{|c|}{\textbf{Method used:} POST} \\
        \hline
        \textbf{Request:} Doctor will post an array of therapeutic segments (e.g., games, study, stories) along with their planned durations. & 
        \textbf{Response:} If status code 201, the session schedule is securely saved and assigned to the child's profile, ready for the robot to fetch and execute. \\
        \hline
    \end{tabular}
\end{table}

\begin{table}[H]
    \centering
    \renewcommand{\arraystretch}{1.5}
    \begin{tabular}{|p{0.46\textwidth}|p{0.46\textwidth}|}
        \hline
        \multicolumn{2}{|c|}{\texttt{/api/posts}} \\
        \hline
        \multicolumn{2}{|c|}{\textbf{API used to publish a new post to the Global Community feed}} \\
        \hline
        \multicolumn{2}{|c|}{\textbf{Method used:} POST} \\
        \hline
        \textbf{Request:} The user (Parent/Doctor) will send the post content, optionally including an image. & 
        \textbf{Response:} If status code 201, the post is published and broadcast to the relevant followers' feeds. \\
        \hline
    \end{tabular}
\end{table}

For testing these APIs, we utilized Postman. It is an API platform for building and using APIs that simplifies each step of the API lifecycle and streamlines collaboration to ensure data integrity before final integration into the mobile client.

\begin{figure}[H]
    \centering
    \includegraphics[width=0.25\textwidth]{Chapter_6/placeholder_postman_testing.png}
    \caption{Postman Logo}
    \label{fig:postman_logo}
\end{figure}

\begin{figure}[H]
    \centering
    \includegraphics[width=\textwidth, height=0.4\textheight, keepaspectratio]{Chapter_6/postman_website.png}
    \caption{API Testing in Postman}
    \label{fig:postman_testing}
\end{figure}

\subsection{Real-Time Communication via Firebase}

While the core relational data of the application (authentication, profiles, the community feed, and session configuration) is served through the Laravel REST API, certain features require persistent, low-latency, two-way communication that traditional request-response cycles are not well suited for. For these specific cases, LUMO integrates Firebase directly into the mobile client, avoiding the overhead of managing a dedicated WebSocket server. The Firebase integration introduces several critical capabilities to the application:

\begin{itemize}
    \item \textbf{Real-Time Chat:} 1-on-1 private messaging between Parent\,$\leftrightarrow$\,Doctor and Doctor\,$\leftrightarrow$\,Doctor is delivered instantly, with messages synchronized live across devices.
    \item \textbf{Typing Indicators:} Efficiently managed and broadcasted across clients utilizing a real-time \texttt{typing\_status} field.
    \item \textbf{Unread Message Counts:} Accurately tracked and synchronized per user through a dedicated \texttt{unread\_counts} map.
    \item \textbf{Push Notifications:} Firebase Cloud Messaging (FCM) delivers instant alerts for new chat messages, new connection requests, and social interactions (likes and comments).
    \item \textbf{Cross-Platform Fallback:} Implements a robust fallback mechanism for Linux Desktop builds. Since Firebase is not natively supported on Linux desktop environments, the application gracefully intercepts these streams and returns \texttt{Stream.value([])}, thereby preventing critical application crashes and ensuring uninterrupted execution.
\end{itemize}

\section{App Services}

The LUMO Mobile Application provides a comprehensive suite of advanced services tailored specifically for both the parents of children with autism and the clinical doctors responsible for their therapy. It is important to note that the mobile application functions as the management and monitoring companion to the LUMO system: the physical robot itself, which directly interacts with the child through voice and motion, is operated through its own dedicated embedded application and is configured remotely via the session schedules built on the mobile app described below.

\subsection{Services for Parents}

\begin{enumerate}
    \item \textbf{Secure Authentication:} Implementation of a robust, role-based registration and authentication system to ensure data privacy and secure access control.
    \item \textbf{Global Community Access:} A centralized social feed enabling parents to share posts, interact with peers, and cultivate broader awareness regarding autism spectrum disorder.
    \item \textbf{AI-Powered Dashboard:} An interactive analytical interface visualizing detailed therapy session outcomes, including quantitative metrics for Focus Scores, Emotion distribution, and Gaze/Eye Contact tracking, once a session has been processed by the AI Cloud Servers.
    \item \textbf{Intelligent Chatbot Integration:} Integration of LUMO's Retrieval-Augmented Generation (RAG) based conversational agent, designed to provide instantaneous, accurate responses to autism-related queries while maintaining a persistent chat history.
    \item \textbf{Real-Time Clinical Communication:} A secure messaging system utilizing Firebase, facilitating real-time communication with assigned doctors, complete with live typing indicators and unread message synchronization.
    \item \textbf{Dynamic Configuration:} Comprehensive profile management and application settings, explicitly supporting Bidi-Aware multi-lingual localization and adaptive theming (Dark Mode).
    \item \textbf{Simplified Progress Tracking:} A dedicated analysis screen showing easy-to-read post-session summaries and trend graphs for their own child, without overwhelming parents with raw clinical data.
\end{enumerate}

\subsection{Services for Doctors}

\begin{enumerate}
    \item \textbf{Verified Clinical Registration:} A stringent registration protocol requiring mandatory identity verification to establish authoritative medical accounts.
    \item \textbf{Patient Management Workflow:} An advanced organizational system allowing practitioners to efficiently receive, review, accept, or reject linkage requests initiated by parents.
    \item \textbf{Centralized Patient Directory:} A consolidated clinical dashboard that shows an aggregate trend overview (indicating Improvement or Stability) for all actively supervised pediatric patients.
    \item \textbf{Therapeutic Session Builder:} A specialized configuration tool enabling doctors to construct bespoke therapy sequences (e.g., Study, Games, Stories, Drawing) and precisely allocate planned temporal durations for each discrete segment. Once saved, the configured session is fetched and executed autonomously by the physical robot.
    \item \textbf{Deep Patient Insights:} Comprehensive data access privileges granting clinicians detailed visibility into the longitudinal analytics and complete historical session data for each linked child.
    \item \textbf{Community Engagement:} Authoritative access to the Global Community, permitting the publication of specialized medical discourse.
    \item \textbf{Remote Real-Time Consultation:} Seamless, low-latency chat functionalities enabling continuous follow-up with parents to provide remote therapeutic guidance and intervention support.
\end{enumerate}
\end{document}
